\section[Ordered path clusters and projective reporters]{Ordered path clusters with persistent projective reporters}
\label{sec:op3-ordered-path-clusters}

\paragraph{Scope.}
The following \statusproved{} statements are proof drafts awaiting independent
review. They give an implemented algebraic backend for supplied rooted tree
clusters. The audit driver rebuilds supplied hierarchies; it does not implement
an online hierarchy or establish a graph-uniform OP3 result. We use exact real
word arithmetic, counting comparisons, unsuccessful searches, dictionary
operations, copies, retained records, graph access and final recovery.
Rational bit complexity and numerical stability are outside this model.
Put $\gamma=(1-\alpha)/(1+\alpha)$, $\lambda=\epsappr/2$,
$\boldsymbol M=\boldsymbol D-\gamma\boldsymbol A$ and
$\boldsymbol b=\boldsymbol e_v-\lambda\boldsymbol d$.
The current positive face $U$ is a connected rooted subtree containing $v$.

\subsection{A finite chart for both endpoint strata}

A two-port cluster has ordered ports $(L,R)$, with $L$ closer to the seed.
Eliminating its interior gives each owned boundary gate the form
\[
 g_j(\boldsymbol u)=a_j u_L+b_j u_R+c_j,
 \qquad a_j,b_j\geq0,\quad a_j+b_j>0,\quad c_j<0.
\]
The source is never an eliminated interior vertex. Consequently the interior
load is negative, and inverse positivity gives the asserted signs.
Use the homogeneous point and its finite chart
\begin{equation}
 (X_j,Y_j,W_j)=(a_j,c_j,a_j+b_j),\qquad
 (x_j,y_j)=(X_j/W_j,Y_j/W_j).
 \label{eq:op3-compact-projective-chart}
\end{equation}
Both $a_j=0$ and $b_j=0$ are finite: $0\leq x_j\leq1$ and $y_j<0$.
The normalized gate is $x_j u_L+(1-x_j)u_R+y_j$. Its maximum is attained
on the upper convex chain of these points. The normalization preserves signs,
so a positive maximizing row is an actual positive original gate, while a
nonpositive maximum certifies every owned row. This does not claim to preserve
the ordering of unnormalized gate magnitudes.

\begin{lemma}[Positive projective pullback; \statusproved]
\label{lem:op3-compact-projective-pullback}
Suppose old ports equal $\boldsymbol P$ times new ports plus
$\boldsymbol t$, where
$\boldsymbol P=\left(\begin{smallmatrix}a&b\\c&d\end{smallmatrix}\right)$
is entrywise nonnegative with positive determinant and
$\boldsymbol t=(t_L,t_R)^\top\leq0$.
The complete row pullback is
\begin{equation}
 \begin{pmatrix}X'\\Y'\\W'\end{pmatrix}=
 \begin{pmatrix}
 a-c&0&c\\ t_L-t_R&1&t_R\\a+b-c-d&0&c+d
 \end{pmatrix}
 \begin{pmatrix}X\\Y\\W\end{pmatrix}.
 \label{eq:op3-compact-projective-matrix}
\end{equation}
It has positive denominator on the whole compact chart, preserves increasing
$x$ order, and maps the old upper convex chain to the new upper convex chain.
\end{lemma}
\begin{proof}
The new coefficients are $aX+c(W-X)$ and $bX+d(W-X)$; their sum is
$W'>0$, including both endpoints since both row sums of $\boldsymbol P$
are positive. Differentiation of $x'=X'/W'$ gives the positive numerator
$ad-bc$. The homogeneous determinant in
\eqref{eq:op3-compact-projective-matrix} is also $ad-bc$.
Every orientation sign is therefore preserved after division by positive
homogeneous coordinates. At fixed $x$, the new $y$ coordinate increases
strictly with the old $y$ coordinate. Thus upper and lower sides cannot
exchange, including equal-$x$ endpoint strata. Negative offsets preserve
$Y'<0$. These facts establish the upper-chain assertion without scanning rows.
\end{proof}

\begin{lemma}[Implemented persistent chain; \statusproved]
\label{lem:op3-persistent-projective-chain}
Let $N$ bound the number of owned rows and put $\ell=\log(2+N)$.
An immutable AVL chain supports a whole-chain pullback in $O(1)$ work and one
copied root; the upper hull of two chains with separated $x$ intervals in
$O(\ell^3)$ work and $O(\ell^2)$ newly allocated nodes; and a gate-sign query
or positive one-port substitution in $O(\ell^2)$ work. Earlier versions remain
queryable. All allocated nodes, including subsequently discarded ones, are
charged to these bounds.
\end{lemma}
\begin{proof}
Store a constant-size homogeneous matrix tag at every node, along with its
point, two children, height and subtree size. A whole-chain map transforms the
root point and composes its tag. A point access carries the product of tags
along one root-to-leaf path in $O(\ell)$ work. Pushing a tag when a structural
operation needs its children copies at most two roots. Persistent AVL split,
join and rotations therefore have at most $O(\ell^2)$ copied nodes and work
under a conservative accounting. The separated-bridge searches from
\cref{lem:op3-avl-upper-hull} use $O(\ell^2)$ point accesses, including failed
comparisons. Homogeneous orientation and endpoint comparisons cost constant
word work per accessed point. After the bridge is found, persistent splits
retain the two surviving subchains and concatenate them. Equal endpoint
abscissae keep the higher point, with a fixed label tie rule; collinear
interior points may be omitted. This proves the stated conservative bounds.

For a one-port substitution
$(u_L,u_R)=(p,q)t+(\xi,\omega)$ with $p,q>0$ and $\xi,\omega\leq0$, a row threshold is
\[
 -\frac{Y+(\xi-\omega)X+\omega W}{(p-q)X+qW}.
\]
The denominator is positive on $[0,1]$. The fractional transformation
$x\mapsto x/((p-q)x+q)$ is increasing, and transforms the upper chain into
an upper chain. The corresponding ordinate, the negative threshold, is thus
unimodal along its ordered vertices. Binary adjacent-value comparisons find
its maximum with $O(\ell)$ point accesses. A flat adjacent maximum is a
supporting edge, so ties cannot conceal a later larger value. Ordinary gate
queries use the same argument with constant denominator. Nodes are immutable;
reclamation, if performed, can be charged once per allocated node. Keeping
all historical versions instead costs at most the total allocation count.
\end{proof}

\subsection{Compression, raking, and one-port reduction}

Cluster edges retain the original diagonal and load at each endpoint.
For a cluster, its matrix uses its own edges and the original diagonal and
load once per vertex. These matrices are positive definite M-matrices:
original degrees dominate the cluster degrees, with strict diagonal dominance
because $0<\gamma<1$. When two edge-disjoint clusters share a port, subtract
one copy of that port's original diagonal and load before eliminating it.
This convention avoids distributing a vertex's load over a changing number
of incident edges.

\begin{lemma}[Adjacent compression has separated chains; \statusproved]
\label{lem:op3-adjacent-chain-separation}
Let two clusters have ports $(L,M)$ and $(M,R)$, with $M\ne v$.
Write their Schur states as $(\boldsymbol H^l,\boldsymbol f^l)$ and
$(\boldsymbol H^r,\boldsymbol f^r)$. Set
\[
 \delta=H^l_{MM}+H^r_{MM}-d_M>0,\quad
 \beta=f^l_M+f^r_M-b_M<0,\quad
 w_L=-H^l_{LM}>0,\quad w_R=-H^r_{MR}>0.
\]
Then
\begin{equation}
 u_M=\eta u_L+\zeta u_R+\sigma,
 \qquad(\eta,\zeta,\sigma)=(w_L,w_R,\beta)/\delta.
 \label{eq:op3-adjacent-recovery}
\end{equation}
The left and right reporter pullbacks have respective matrices
$\left(\begin{smallmatrix}1&0\\\eta&\zeta\end{smallmatrix}\right)$ and
$\left(\begin{smallmatrix}\eta&\zeta\\0&1\end{smallmatrix}\right)$,
and offsets $(0,\sigma)$ and $(\sigma,0)$. Their new compact abscissae lie in
\[
 \text{right child: }[0,\eta/(\eta+\zeta)],\qquad
 \text{left child: }[\eta/(\eta+\zeta),1].
\]
Hence their exact union reporter uses one separated merge, with the right
child first. The new Schur matrix, load and recovery record need $O(1)$ words.
\end{lemma}
\begin{proof}
Schur elimination gives \eqref{eq:op3-adjacent-recovery}. Positivity follows
from the M-matrix property; $\beta<0$ follows because every eliminated vertex
is a nonseed with negative original load. Each displayed port matrix is
nonnegative with positive determinant, so
\cref{lem:op3-compact-projective-pullback} applies. Evaluate the monotone
abscissa maps at $0$ and $1$ to obtain the intervals. They meet only at the
shared-port endpoint. The new entries are
\[
 H_{LL}=H^l_{LL}-w_L\eta,\quad
 H_{RR}=H^r_{RR}-w_R\zeta,\quad
 H_{LR}=-w_L\zeta=-w_R\eta,
\]
with loads $f_L=f^l_L+w_L\sigma$ and $f_R=f^r_R+w_R\sigma$.
\end{proof}

\paragraph{Raking and forgetting.}
A one-port side cluster at $i$ is represented by its scalar Schur state and
its smallest boundary threshold $\theta$. Raking it into a path cluster adds
its diagonal and load, subtracting the original $d_i,b_i$ once. Conditional
on the unchanged path ports, the old path interior and reporter do not
change. All side gates reduce exactly to the single inequality $u_i>\theta$,
so its reporter contributes one endpoint point, merged at $x=0$ or $x=1$.
Raking two one-port clusters takes the minimum threshold and the same scalar
state sum. If there are no owned gates, omit the threshold.

To forget the farther port $R\ne v$ of a two-port cluster, substitute
\[
 u_R=(-H_{LR}/H_{RR})u_L+f_R/H_{RR}.
\]
Its multiplier is positive and offset negative. Apply the one-port query of
\cref{lem:op3-persistent-projective-chain}, update the scalar Schur state,
and save the affine recovery record. This proves closure of all these
operations with $O(\ell^3)$ work and $O(\ell^2)$ allocated hull nodes per
operation, independently of how many original rows the cluster represents.

\paragraph{Unique row ownership and recovery.}
Every inactive boundary vertex $j$ of a connected active subtree has exactly
one active parent $i$. At $i$, maintain the minimum of the local thresholds
$\lambda d_j/\gamma$ in a heap. All these gates have the same coefficient
of $u_i$, so this single minimum is an exact local certificate. Home the
resulting single row on $i$'s permanent parent edge; home the seed's row on
its first admitted edge. An admission changes the old parent's one home-edge
payload and creates the new vertex's payload. Copying that row to every
incident edge would be unnecessary and could cost quadratically on a star.
The singleton face is handled directly before the first edge exists.
Original adjacency entries and neighbor-degree queries are charged when a
vertex is admitted, with cached records and heap construction also charged.

The current hierarchy has a linear number of cluster records. A final
traversal uses the saved affine recovery equations, passes port values to
children, and writes each active vertex once. A dictionary verifies or
coalesces repeated shared ports; balanced dictionaries give $O(|U|\log(2+|U|))$
work, or a preassigned active index gives linear word work. Historical roots
need not be traversed. The returned PageRank vector is
$\bar\alpha\boldsymbol D\boldsymbol u$, with original residual
$\boldsymbol e_v-\boldsymbol M\boldsymbol u$.

\subsection{Precise remaining hierarchy contract}

\paragraph{Source.}
Alstrup et al.\ \cite{alstrup2005toptrees}, journal Section 2 and
Theorem 2.1 (pp.\ 245--248), maintain edge-induced clusters with at
most two boundary vertices. A link, cut or expose uses $O(\log n)$ joins and
splits, with logarithmic height and linear structural space. Section 6.1
(pp.\ 259--260) explains exposure; Section 6.2 (p.\ 260) handles arbitrary
degrees within the implementation. Link and cut reset external boundaries.
This source supplies a hierarchy theorem, not the PageRank summaries above.

\paragraph{Conditional.}
If an online hierarchy maintains the stated rooted cluster summaries with
$O(\log(2+|U|))$ changed records per admission and home-edge update, then the
backend yields total word work
\[
 O\bigl(\operatorname{cvol}(U)\log^4(2+\operatorname{cvol}(U))\bigr)
\]
and at most $O(\operatorname{cvol}(U)\log^3(2+\operatorname{cvol}(U)))$
retained state, even keeping every old application version. Temporary
query matrices are discarded as the search advances; their allocations
are charged to word work, not accumulated as persistent hull nodes. The root's one-port
Schur equation gives $u_v$ exactly, and its minimum threshold gives either a
strict positive original gate or a complete quiet certificate. Every strict
admission increases the positive face, whose volume is at most
$1/\lambda=2/\epsappr$. Thus this contract would give
$\widetilde O(1/\epsappr)$ local ACL work on arbitrary trees. This statement
charges failed searches, the final quiet query, local heaps, graph discovery,
all factor and hull records, dictionaries and output.

\paragraph{Integration obligations at the primitive stage.}
The code currently supplies and reparents static hierarchies. Before removing
the preceding condition, reconcile every join shape, singleton initialization,
source exposure during updates, changing base-edge payloads and dynamic
vertex allocation with the source interface. In particular, a temporary
cluster with an interior seed cannot use the negative-load argument. One
candidate is to mark such a summary unavailable until seed exposure splits
it; another is the source's composite-update interface. The unavailable-summary option is proved and fully integrated in
\cref{sec:op3-local-top-trees}; it is not silently assumed by this primitive. The source's internal high-degree representation is
structural and must not become extra PageRank vertices. No ambient graph
preprocessing is permitted.

\paragraph{Measured.}
The implemented immutable chain passes 516 independent materialized-chain
comparisons, 360 exact extreme queries, 360 one-port queries and 96 adjacent
separation checks, with 130 retained versions rechecked. The cluster audit
uses every nonisomorphic tree through seven vertices, every seed and three
parameter pairs, following exact positive faces. It checks independent dense
Schur complements, every raw boundary inequality, recovered coordinates,
three path parenthesizations and both orders of one-port rakes. Five implicit
unweighted path/private-star families through 128 active vertices keep all
boundary rows on the upper chain; their inactive adjacency is never scanned.
These are correctness and representation audits, not an online hierarchy
benchmark. Exact counts and source hashes are in
\texttt{PROJECTIVE\_HULL\_ROPE\_AUDIT.json} and
\texttt{PATH\_CLUSTER\_REPORTER\_AUDIT.json}.
